% GENERATED FILE. Edit canonical lessons, then run npm run generate:books.
% canonical-source-hash: fnv1a64:511e5b88e307b146
% canonical-lessons: AR-C248-intarnit, AR-C248-mawqi, AR-C248-tatbiq, AR-C248-jihaz, AR-C248-musalsal

\chapter{Internet, Website, App, Device, TV series}
\label{ch:ar-internet-website-app-device-tv-series}
\hlchaptermodality{248}

\begin{chapteropening}
\textbf{By the end of this chapter:} \emph{I can say the internet, a website, an app, a device and a TV series.}

Complete the last lesson of chapter 248: I can say the internet, a website, an app, a device and a TV series.
\end{chapteropening}

\section[\texorpdfstring{\ar{إنترنت} (intarnit)}{إنترنت (intarnit)}]{\ar{إنترنت} (intarnit) — the internet}

\label{lesson:AR-C248-intarnit}



\begin{quote}
Before the new one: say the Arabic for a palm tree, then the Arabic for jasmine.
\end{quote}

\subsection*{You'll want to know: \ar{إنترنت}}
\textbf{\ar{إنترنت}} — \emph{intarnit} — \textquotedblleft{}the internet\textquotedblright{}.

\begin{cousinweb}
\textbf{\ar{إنترنت}} (\emph{intarnit}) is a modern loan from English \emph{internet} (EtymDB); it has no Arabic root.
\end{cousinweb}

\begin{grammarlens}[title={What you've built}]
One more word for this chapter.
\end{grammarlens}

\subsection*{Your turn}
\begin{itemize}
\raggedright
  \item \emph{Say it:} \emph{intarnit}
  \item \emph{Say it:} \emph{intarnit}, once more
  \item \emph{Say it:} say \emph{intarnit}
  \item \emph{Recall:} say the Arabic for a palm tree, then the Arabic for jasmine, then say \emph{intarnit} again
\end{itemize}

\begin{tcolorbox}[breakable,colback=teal!4,colframe=teal!35!black,title={Before you move on}]
What does \ar{إنترنت} mean? (\textquotedblleft{}The internet\textquotedblright{}.) Say it once more.
\end{tcolorbox}
\section[\texorpdfstring{\ar{موقع} (mawqiʿ)}{موقع (mawqiʿ)}]{\ar{موقع} (mawqiʿ) — a website, a site}

\label{lesson:AR-C248-mawqi}



\begin{quote}
Before the new one: say the Arabic for jasmine, then the Arabic for the internet.
\end{quote}

\subsection*{You'll want to know: \ar{موقع}}
\textbf{\ar{موقع}} — \emph{mawqiʿ} — \textquotedblleft{}a website, a site\textquotedblright{}.

\begin{cousinweb}
\textbf{\ar{موقع}} (\emph{mawqiʿ}) is built on the native Arabic root \emph{w-q-ʿ}, \textquotedblleft{}to fall\textquotedblright{} (Lane). \emph{Mawqiʿ} is the \emph{ma-} place noun, \textquotedblleft{}a place where a thing falls\textquotedblright{}, hence \textquotedblleft{}a site\textquotedblright{}; the website is a modern sense.
\end{cousinweb}

\begin{grammarlens}[title={What you've built}]
One more word for this chapter.
\end{grammarlens}

\subsection*{Your turn}
\begin{itemize}
\raggedright
  \item \emph{Say it:} \emph{mawqiʿ}
  \item \emph{Say it:} \emph{mawqiʿ}, once more
  \item \emph{Say it:} say \emph{mawqiʿ}
  \item \emph{Recall:} say the Arabic for jasmine, then the Arabic for the internet, then say \emph{mawqiʿ} again
\end{itemize}

\begin{tcolorbox}[breakable,colback=teal!4,colframe=teal!35!black,title={Before you move on}]
What does \ar{موقع} mean? (\textquotedblleft{}A website, a site\textquotedblright{}.) Say it once more.
\end{tcolorbox}
\section[\texorpdfstring{\ar{تطبيق} (taṭbīq)}{تطبيق (taṭbīq)}]{\ar{تطبيق} (taṭbīq) — an app, an application}

\label{lesson:AR-C248-tatbiq}



\begin{quote}
Before the new one: say the Arabic for the internet, then the Arabic for a website.
\end{quote}

\subsection*{You'll want to know: \ar{تطبيق}}
\textbf{\ar{تطبيق}} — \emph{taṭbīq} — \textquotedblleft{}an app, an application\textquotedblright{}.

\begin{cousinweb}
\textbf{\ar{تطبيق}} (\emph{taṭbīq}) is built on the native Arabic root \emph{ṭ-b-q}, \textquotedblleft{}to cover, to fit one on another\textquotedblright{} (Lane). \emph{Taṭbīq} is the verbal noun of Form II, \textquotedblleft{}applying\textquotedblright{}; the app is a modern sense.
\end{cousinweb}

\begin{grammarlens}[title={What you've built}]
One more word for this chapter.
\end{grammarlens}

\subsection*{Your turn}
\begin{itemize}
\raggedright
  \item \emph{Say it:} \emph{taṭbīq}
  \item \emph{Say it:} \emph{taṭbīq}, once more
  \item \emph{Say it:} say \emph{taṭbīq}
  \item \emph{Recall:} say the Arabic for the internet, then the Arabic for a website, then say \emph{taṭbīq} again
\end{itemize}

\begin{tcolorbox}[breakable,colback=teal!4,colframe=teal!35!black,title={Before you move on}]
What does \ar{تطبيق} mean? (\textquotedblleft{}An app, an application\textquotedblright{}.) Say it once more.
\end{tcolorbox}
\section[\texorpdfstring{\ar{جهاز} (jihāz)}{جهاز (jihāz)}]{\ar{جهاز} (jihāz) — a device}

\label{lesson:AR-C248-jihaz}



\begin{quote}
Before the new one: say the Arabic for a website, then the Arabic for an app.
\end{quote}

\subsection*{You'll want to know: \ar{جهاز}}
\textbf{\ar{جهاز}} — \emph{jihāz} — \textquotedblleft{}a device\textquotedblright{}.

\begin{cousinweb}
\textbf{\ar{جهاز}} (\emph{jihāz}) is built on the native Arabic root \emph{j-h-z}, \textquotedblleft{}to equip, to make ready\textquotedblright{} (Lane). \emph{Jihāz} is equipment, an apparatus.
\end{cousinweb}

\begin{grammarlens}[title={What you've built}]
One more word for this chapter.
\end{grammarlens}

\subsection*{Your turn}
\begin{itemize}
\raggedright
  \item \emph{Say it:} \emph{jihāz}
  \item \emph{Say it:} \emph{jihāz}, once more
  \item \emph{Say it:} say \emph{jihāz}
  \item \emph{Recall:} say the Arabic for a website, then the Arabic for an app, then say \emph{jihāz} again
\end{itemize}

\begin{tcolorbox}[breakable,colback=teal!4,colframe=teal!35!black,title={Before you move on}]
What does \ar{جهاز} mean? (\textquotedblleft{}A device\textquotedblright{}.) Say it once more.
\end{tcolorbox}
\section[\texorpdfstring{\ar{مسلسل} (musalsal)}{مسلسل (musalsal)}]{\ar{مسلسل} (musalsal) — a TV series}

\label{lesson:AR-C248-musalsal}



\begin{quote}
Before the new one: say the Arabic for an app, then the Arabic for a device.
\end{quote}

\subsection*{You'll want to know: \ar{مسلسل}}
\textbf{\ar{مسلسل}} — \emph{musalsal} — \textquotedblleft{}a TV series\textquotedblright{}.

\begin{cousinweb}
\textbf{\ar{مسلسل}} (\emph{musalsal}) is the passive participle of \emph{salsala}, \textquotedblleft{}to link in a chain\textquotedblright{}, built on the doubled root \emph{s-l-s-l}, which Lane lists under \emph{s-l-l} with \emph{silsila}, \textquotedblleft{}a chain\textquotedblright{}: linked one to the next, a serial.
\end{cousinweb}

\begin{grammarlens}[title={What you've built}]
One more word for this chapter.
\end{grammarlens}

\subsection*{Your turn}
\begin{itemize}
\raggedright
  \item \emph{Say it:} \emph{musalsal}
  \item \emph{Say it:} \emph{musalsal}, once more
  \item \emph{Say it:} say \emph{musalsal}
  \item \emph{Recall:} say the Arabic for an app, then the Arabic for a device, then say \emph{musalsal} again
\end{itemize}

\begin{tcolorbox}[breakable,colback=teal!4,colframe=teal!35!black,title={Before you move on}]
What does \ar{مسلسل} mean? (\textquotedblleft{}A TV series\textquotedblright{}.) Say it once more.
\end{tcolorbox}
